\section{Diseño experimental y métricas}

\subsection{Validación sintética}

Para cada señal sintética, el segundo bloque se mantuvo como conjunto de prueba. Los tres modelos recibieron la misma partición. Se evaluaron RMSE, MAE, $R^2$ y correlación de Pearson de la señal retenida. Para la HRF se utilizaron RMSE, $R^2$, correlación, error de amplitud máxima y error en tiempo al pico. En la PINN también se calcularon errores relativos de $\epsilon$, $\tau$ y $\alpha$.

El diseño pareado permitió aplicar una prueba de Friedman por nivel de ruido y comparaciones de Wilcoxon con corrección de Holm. Los tamaños de efecto se expresaron mediante correlación biserial por rangos.

\subsection{Evaluación real bidireccional}

Cada escenario real contiene dos bloques. Se realizaron dos direcciones:

\begin{enumerate}
    \item entrenar en bloque 1 y evaluar en bloque 2;
    \item entrenar en bloque 2 y evaluar en bloque 1.
\end{enumerate}

Con cuatro escenarios se obtuvieron ocho experimentos por modelo. La métrica primaria fue RMSE fuera de muestra. Además, se estudió la brecha de generalización definida como RMSE de prueba menos RMSE de entrenamiento.

\subsection{Consistencia fisiológica}

Para cada escenario se compararon las HRF y parámetros estimados desde bloques distintos. La consistencia de HRF se cuantificó mediante RMSE y correlación de Pearson. Para los parámetros se utilizó la diferencia porcentual simétrica.
